\documentclass[11pt,a4paper]{article}
\usepackage[T1]{fontenc}
\usepackage[utf8]{inputenc}
\usepackage[french]{babel}
\frenchsetup{SmallCapsFigTabCaptions=false}
\usepackage{lmodern,microtype}
\usepackage[left=22mm,right=22mm,top=23mm,bottom=24mm]{geometry}
\usepackage{graphicx,booktabs,array,tabularx,longtable,float}
\usepackage{xcolor,tikz}
\usepackage{enumitem,titlesec}
\usepackage{hyperref}
\hypersetup{colorlinks=true,linkcolor=black,citecolor=black,urlcolor=sfigBlue,pdfauthor={Projet LANCE},pdflang={fr-FR}}
\urlstyle{same}
\setlength{\parindent}{1.2em}
\setlength{\parskip}{3pt}
\setlength{\emergencystretch}{2em}
\clubpenalty=10000
\widowpenalty=10000
\displaywidowpenalty=10000
\setlist{nosep,leftmargin=5mm,topsep=3pt}
\def\UrlBreaks{\do\/\do\-\do\_\do\.\do\?\do\&\do\=\do\:\do\,}

\newcolumntype{Y}{>{\raggedright\arraybackslash}X}
\newcommand{\code}[1]{{\small\ttfamily\path{#1}}}
\newcommand{\evidence}[1]{\par\smallskip{\small\color{sfigMuted}\textbf{Preuve / limite.} #1}\par\smallskip}
\titleformat{\section}{\Large\bfseries}{\thesection}{0.7em}{}
\titleformat{\subsection}{\large\bfseries}{\thesubsection}{0.7em}{}
\titlespacing*{\section}{0pt}{16pt}{7pt}
\titlespacing*{\subsection}{0pt}{11pt}{5pt}
\pagestyle{plain}
\setcounter{tocdepth}{1}

% Reusable TikZ styles for sober scientific figures.
% Required packages: xcolor, tikz.
% Required TikZ libraries: arrows.meta, positioning, fit, calc.

\usetikzlibrary{arrows.meta,positioning,fit,calc}

\definecolor{sfigInk}{HTML}{1F242C}
\definecolor{sfigMuted}{HTML}{5C626C}
\definecolor{sfigLine}{HTML}{D9DEE5}
\definecolor{sfigNeutral}{HTML}{F7F8FA}
\definecolor{sfigBlue}{HTML}{1D5C96}
\definecolor{sfigBlueLine}{HTML}{93B1CD}
\definecolor{sfigBlueFill}{HTML}{EFF6FC}
\definecolor{sfigRaspberry}{HTML}{AF0032}
\definecolor{sfigRoseLine}{HTML}{DB8CA3}
\definecolor{sfigRoseFill}{HTML}{FDF0F4}

\tikzset{
  sfig diagram/.style={
    >=Latex,
    node distance=10mm and 12mm,
    every node/.style={font=\sffamily\small,text=sfigInk}
  },
  sfig base/.style={
    draw=sfigLine,
    fill=white,
    rounded corners=2pt,
    line width=0.55pt,
    inner xsep=6pt,
    inner ysep=4pt,
    align=center
  },
  sfig input/.style={
    sfig base,
    draw=sfigBlueLine,
    fill=sfigBlueFill
  },
  sfig process/.style={
    sfig base,
    draw=sfigRoseLine,
    fill=sfigRoseFill
  },
  sfig neutral/.style={
    sfig base,
    fill=sfigNeutral
  },
  sfig focal/.style={
    sfig base,
    draw=sfigRaspberry,
    fill=sfigRoseFill,
    line width=0.75pt,
    font=\sffamily\small\bfseries,
    text=sfigRaspberry
  },
  sfig output/.style={
    sfig neutral,
    font=\sffamily\small\bfseries
  },
  sfig group/.style={
    draw=sfigLine,
    fill=none,
    dashed,
    rounded corners=3pt,
    line width=0.5pt,
    inner sep=6pt
  },
  sfig flow/.style={
    -{Latex[length=2mm,width=1.2mm]},
    draw=sfigMuted,
    line width=0.6pt
  },
  sfig blue flow/.style={
    -{Latex[length=2mm,width=1.2mm]},
    draw=sfigBlue,
    line width=0.68pt
  },
  sfig primary flow/.style={
    -{Latex[length=2mm,width=1.2mm]},
    draw=sfigRaspberry,
    line width=0.78pt
  },
  sfig blocked/.style={
    draw=sfigRaspberry,
    dashed,
    line width=0.68pt
  },
  sfig annotation/.style={
    font=\sffamily\scriptsize,
    text=sfigMuted,
    align=center
  },
  sfig formula/.style={
    font=\small,
    text=sfigInk,
    align=center
  }
}

% Classical article front matter; standalone source for the native editor.
\renewenvironment{abstract}
  {\begin{quote}\small\noindent\textbf{Résumé.}\ }
  {\par\end{quote}}
\title{\bfseries LANCE : proposition de spécialisation IoMT\\[0.5em]\large\mdseries\itshape Périmètre simulé, vérification des preuves et évaluation comparative}
\author{Projet LANCE}
\date{Dernière révision : 1er octobre 2026 à 10:23\\[0.15em]\small Europe/Brussels (UTC+02:00)\\[0.4em]\small Proposition de TFE ; aucun résultat expérimental}
\hypersetup{pdftitle={LANCE : proposition de spécialisation IoMT}}
\begin{document}
\maketitle
\begin{abstract}
Cette proposition de travail de fin d'études (TFE) examine comment adapter LANCE,
un outil d'audit autorisé de l'Internet des objets (IoT), à des services médicaux
simulés de l'Internet of Medical Things (IoMT). La relecture de la proposition
historique fait apparaître trois points à préciser : services effectivement
couverts, preuves attendues et portée des conclusions. Le travail proposé
consiste à construire des scénarios vulnérables et durcis, puis à comparer la
version généraliste et la version spécialisée à ressources comparables.
Les observations passives et les actions actives de laboratoire sont évaluées
séparément. Un éventuel bénéfice technique reste à mesurer ; le banc virtuel
ne valide ni la sûreté d'un dispositif médical réel ni un déploiement hospitalier.
\end{abstract}
\clearpage
\begingroup\small\tableofcontents\endgroup
\clearpage

\section{Question et périmètre}
La question est de savoir si une spécialisation IoMT améliore la découverte des
services et la confirmation des propriétés de sécurité sur un banc simulé,
par rapport à la version généraliste de LANCE. L'amélioration est une hypothèse,
pas une conséquence garantie de l'ajout de connaissances médicales.

Le périmètre envisagé comprend des services d'imagerie, de messages cliniques
et de télémétrie, alimentés par des données synthétiques. Il ne comprend aucun
patient, dispositif de soin réel ou donnée de santé réelle. Les essais futurs
restent à autoriser explicitement et ciblent nato/pve-nato selon le guide du
dépôt~\cite{repo}. Aucune infrastructure n'a été modifiée pour cette relecture.

Un grand modèle de langage (LLM) peut participer aux décisions de LANCE.
La présente comparaison porte sur la spécialisation du système complet ; elle
n'isole pas automatiquement l'apport causal du modèle. Pour cette autre question,
le protocole agent/automatisation compare D1, politique à règles sans LLM,
et A1, agent avec LLM, dans un périmètre borné~\cite{agent}.

\section{Méthode de la relecture et provenance}
La date de rédaction de la proposition initiale est inconnue. Le dossier a été
classé le 29 septembre 2026 ; cette date ne constitue pas une date d'expérience.
La présente révision relit \code{subject.md}, les conventions de recherche et
les pages de référence sur les scénarios et l'évaluation. La base du checkout
est \code{5f78e5c9eba4dbbe51ce91c08176a50b611928e1}, avec des modifications
documentaires locales ; ce SHA ne représente pas à lui seul les textes révisés.

Le document Word \code{subject.docx} est conservé comme version historique.
Son contenu n'a pas été rendu ni corrigé lors de cette révision. Le présent
rapport et sa source LaTeX portent le cadre révisé. Le README sert de résumé
et de navigation ; \code{subject.md} et \code{subject.docx} conservent la
proposition historique. Aucun code IoMT, service réseau, rôle Ansible ou score
n'a été validé par cette relecture.

Deux sources primaires ont été consultées le 1er octobre 2026 : la section de
sécurité DICOM PS3.15 B.12, version 2026d~\cite{dicom}, et les sections HTTP et
DICOM de la documentation de sécurité Orthanc~\cite{orthanc}. Elles servent à
préciser ce qui devra être testé. Cette lecture ciblée n'est ni une revue
exhaustive des normes médicales ni une validation clinique.

\section{Constats et points à clarifier}
La proposition historique mélangeait la vérification dite non intrusive et
un chemin comprenant compromission et pivot. Une observation passive lit une
trace déjà disponible ; une requête active sollicite un service, même si elle
ne vise pas à le modifier. Une action de compromission ou une modification de
service appartient à une autre classe d'essai. Leur présence dans un banc
simulé ne prouve pas leur innocuité sur un appareil médical.

La présence d'un protocole n'est pas, à elle seule, une faille. DICOM décrit un
profil de transport sécurisé par TLS~\cite{dicom}. Orthanc distingue la sécurité
de son interface HTTP et celle de son serveur DICOM~\cite{orthanc}. Le scénario
doit donc préciser interface, identité, opération et résultat attendu avant de
qualifier une exposition. Une simple bannière de service ne prouve pas un défaut
d'authentification, ni une fuite de données.

La pondération clinique envisagée n'est pas le score principal actuel de LANCE.
Ce dernier repose sur le F1 final des confirmations ; les scores de sévérité
restent distincts~\cite{evaluation}. Une nouvelle grille clinique serait un
instrument de recherche à construire et à faire valider, sans modifier
implicitement ce contrat logiciel.

\section{Services et propriétés proposés}
Le tableau~\ref{tab:scope} distingue les rôles applicatifs et les limites du
banc envisagé. Il ne décrit pas des extensions déjà disponibles.
PACS désigne un système d'archivage et de communication des images ; DICOM le
standard utilisé pour les données d'imagerie. HL7 désigne une famille de standards
d'échange de données cliniques. MQTT est un protocole de publication/abonnement ;
BLE désigne Bluetooth Low Energy.

\begin{table}[H]\centering\small
\begin{tabularx}{\linewidth}{@{}>{\raggedright\arraybackslash}p{31mm}YY@{}}\toprule
Rôle proposé & Propriété à tester & Limite de représentation \\\midrule
Imagerie PACS / DICOM & Lecture ou transfert autorisé/refusé, chiffrement de l'interface effectivement utilisée. & Orthanc ne représente pas tous les équipements d'imagerie ; distinguer DICOM et API HTTP. \\
Messages HL7 & Règles d'accès et de traitement d'un profil et d'un transport fixés avant test. & Version des messages et transport à choisir ; un parseur permissif ne suffit pas à définir une injection. \\
Télémétrie MQTT & Accès à un topic synthétique selon l'identité, lecture autorisée et refus témoin. & Aucun moniteur clinique réel ni preuve de fidélité physiologique. \\
Moniteur ou pompe émulés & Authentification et propriété applicative explicitement simulée. & Une API HTTP ne représente pas la sûreté physique d'une pompe ni son firmware réel. \\
BLE & Hors du premier périmètre proposé. & La couche radio du banc envisagé reste non validée ; une émulation TCP ne prouverait pas sa couverture. \\\bottomrule
\end{tabularx}
\caption{Périmètre envisagé : chaque propriété exige un attendu positif, un contrôle négatif et une preuve propre à l'interface.}\label{tab:scope}
\end{table}

Pour chaque cas, la vérité terrain (\emph{ground truth}, GT) décrit la propriété
attendue et reste réservée à l'évaluateur. Un contrôle négatif exige un
comportement sain précis, par exemple refuser la lecture d'un topic à une
identité limitée. Conserver des cas entièrement durcis pour mesurer les
fausses alertes ; un scénario contenant quelques contrôles sains ne suffit pas.

Le chemin poste soignant--télémétrie--PACS reste une \emph{hypothèse de scénario}.
La topologie, les accès et les dépendances doivent le rendre possible ; aucune
relation réseau ne se déduit des seuls rôles médicaux. Les expériences de pivot
sont facultatives et distinctes du premier audit de services. Elles exigeraient
une preuve de transition depuis une session source et un contrôle d'absence
d'accès direct, conformément aux limites déjà étudiées~\cite{benchmark}.

\section{Protocole d'évaluation proposé}
La figure~\ref{fig:protocol} présente l'ordre de travail. Chaque étape reste
proposée, y compris la construction et le déploiement des services.

\begin{figure}[H]\centering
\begin{tikzpicture}[sfig diagram]
\node[sfig input,text width=27mm,minimum height=15mm] (scope) at (0,0) {Propriétés\\et preuves};
\node[sfig base,text width=28mm,minimum height=15mm] (cases) at (4.1,0) {Cas vulnérables\\et durcis};
\node[sfig focal,text width=28mm,minimum height=15mm] (pair) at (8.2,0) {Audits\\appariés};
\node[sfig output,text width=28mm,minimum height=15mm] (results) at (12.3,0) {Résultats\\et limites};
\draw[sfig flow] (scope) -- (cases);
\draw[sfig primary flow] (cases) -- (pair);
\draw[sfig flow] (pair) -- (results);
\end{tikzpicture}
\caption{Évaluation proposée : comparer les deux versions sur les mêmes états initiaux, après validation des propriétés et des preuves. Aucun gain n'est présumé.}\label{fig:protocol}
\end{figure}

\subsection{Comparaison principale}
Figer une version généraliste et une version spécialisée, puis les exécuter sur
les mêmes instances et états initiaux restaurés. Conserver le même modèle,
les mêmes paramètres et des plafonds comparables de durée et de charge réseau.
Documenter les outils et connaissances ajoutés : ce contraste mesure leur effet
combiné. Une comparaison supplémentaire où seuls les prompts changent pourrait
isoler leur contribution, si les mêmes observations et outils restent accessibles.

Séparer les cas utilisés pour le développement de variantes réservées avant
les réglages. Définir les répétitions à partir d'un pilote et de sa variance ;
aucun nombre d'essais n'est justifié à ce stade. Enregistrer échecs, exclusions,
états incomplets et coûts, avec une règle de disponibilité fixée avant lecture
des résultats. Une mesure manquante ne devient pas un zéro.

\subsection{Mesures et qualification des preuves}
Un vrai positif (VP) est une propriété attendue correctement confirmée avec
la preuve exigée. Pour $N$ prédictions dédupliquées et $G>0$ propriétés positives
attendues, le contrat actuel donne précision $VP/N$ si $N>0$, rappel $VP/G$ et
F1 $2VP/(N+G)$~\cite{evaluation}. Le F1 combine la précision et le rappel.
Sur les cas entièrement sains, analyser séparément les fausses alertes : leur
rappel et leur F1 restent indisponibles dans ce contrat.

Mesurer aussi couverture des services, taux de preuves acceptables, durée,
charge de sondage et coût total, y compris les essais échoués. Les observations
passives, requêtes actives de lecture et actions modifiant un service sont
rapportées séparément. Une grille d'impact clinique serait secondaire,
préenregistrée et revue avec des personnes qualifiées ; aucun score clinique
n'est retenu dans cette proposition.

Chaque propriété simulée exige un validateur et des contre-exemples : réponse
d'aide, observation d'une mauvaise identité, ou simple déclaration de succès
ne doivent pas recevoir automatiquement le crédit d'une confirmation.
La revue des traces doit être indépendante de la politique comparée.

\subsection{Conditions de reproduction et livrables}
Avant toute expérience, conserver code et modifications locales, images et
versions de services, topologie, données synthétiques, identités de test,
contrats de preuve, configuration des modèles, budgets, ordre des essais et
attestation de remise à zéro. Référencer ensuite les identifiants de runs et
leurs artefacts, sans publier de secret. Aucun de ces artefacts expérimentaux
n'existe à ce stade pour la spécialisation proposée.

Les livrables envisagés sont une extension logicielle, des scénarios
reproductibles, une analyse comparative et un rapport de limites. Leur
validation logicielle et leur évaluation expérimentale seront rapportées
séparément. Une proposition déclarative de scénario ne garantit ni son
déploiement ni la présence des propriétés attendues.

\section{Discussion et limites}
Le banc virtuel permettrait d'évaluer des propriétés réseau et applicatives
explicitement simulées. Il ne reproduit pas automatiquement comportement radio,
firmware, matériel, contraintes temporelles cliniques ou effets sur un patient.
Un gain sur ces services ne permettrait pas d'annoncer la sûreté d'un audit en
soins, une conformité réglementaire ou l'aptitude à un déploiement hospitalier.

Une spécialisation peut améliorer la couverture grâce à de nouveaux outils,
à des règles ou à des connaissances de domaine, sans que le LLM soit la cause du
gain. Les effets doivent être attribués au contraste effectivement réalisé.
Des coûts ou fausses alertes accrus peuvent rendre le résultat défavorable ;
une conclusion négative ou indéterminée reste admissible.

La faisabilité des rôles médicaux, la définition du transport HL7, la capacité
de durcissement de chaque service, les preuves de pivot et les pondérations
cliniques restent à vérifier. Cette relecture documentaire ne les résout pas.

\section{Conclusion et prochaines étapes}
Le sujet peut être formulé comme une étude comparative de spécialisation sur
services IoMT simulés. Les étapes prioritaires sont de choisir un petit ensemble
de propriétés vérifiables, définir les preuves et les témoins durcis, confirmer
la faisabilité du banc puis fixer le protocole avant expérimentation.
L'extension BLE et les pivots doivent être étudiés séparément si leur utilité
est justifiée. Les conclusions finales porteront sur les services testés et
leurs limites ; les recommandations hospitalières resteront préliminaires,
à valider dans un cadre distinct.

\begin{thebibliography}{9}
\phantomsection\addcontentsline{toc}{section}{Références}
\interlinepenalty=10000
\small
\bibitem{repo} Projet LANCE. Guide du dépôt, conventions de recherche et description du pipeline. Documents locaux consultés le 1er octobre 2026, base de checkout \code{5f78e5c9eba4dbbe51ce91c08176a50b611928e1} avec modifications documentaires locales. \code{AGENTS.md}, \code{research/README.md}, \code{src/agent/phases/README.md}. La proposition IoMT reste une piste de recherche.
\bibitem{evaluation} Projet LANCE. \emph{Évaluation et preuves V1}. Sections sur populations, métriques et qualification de preuve lues le 1er octobre 2026. \code{docs/benchmark/v1/evaluation.md}. Aucun score IoMT reproduit.
\bibitem{agent} Projet LANCE. \emph{LANCE : mesurer l'apport des LLM}. Étude commencée le 27 septembre 2026 ; revue des constats du 29 septembre. \href{../2026-09-27-agent-vs-automation/report.pdf}{Rapport}. Protocole proposé, gain expérimental non établi.
\bibitem{benchmark} Projet LANCE. \emph{Couverture et validité du benchmark}. Étude du 29 septembre 2026, limites de preuves de chemins et préparation des propriétés. \href{../2026-09-29-benchmark-coverage/report.pdf}{Rapport}. Aucun résultat clinique.
\bibitem{dicom} DICOM Standards Committee. \emph{DICOM PS3.15 2026d, B.12 : BCP 195 RFC 8996, 9325 TLS Secure Transport Connection Profile}. Section B.12 lue le 1er octobre 2026 ; périmètre limité au profil de transport, pas à toute la norme. \url{https://dicom.nema.org/medical/dicom/current/output/chtml/part15/sect_B.12.html}. Le lien courant est évolutif.
\bibitem{orthanc} Communauté Orthanc. \emph{Securing Orthanc}. Documentation vivante, date de publication non précisée ; sections de sécurité HTTP et DICOM lues le 1er octobre 2026. \url{https://orthanc.uclouvain.be/book/faq/security.html}. Aucun déploiement reproduit.
\end{thebibliography}
\end{document}
